\documentclass[11pt,a4paper]{article}

\usepackage[margin=1in]{geometry}
\usepackage{amsmath,amssymb,amsthm}
\usepackage{enumitem}
\usepackage{graphicx}
\usepackage{hyperref}
\usepackage{url}

\hypersetup{
  colorlinks=true,
  linkcolor=blue,
  citecolor=blue,
  urlcolor=blue
}

\newtheorem{theorem}{Theorem}[section]
\newtheorem{lemma}[theorem]{Lemma}

\newcommand{\C}{\mathbb{C}}
\newcommand{\Arg}{\operatorname{arg}}

\title{Fully Radical Solution of the Quartic Equation over \(\C\)}
\author{LiZhi Zhou\\[0.3em]\small Project artifact: \url{https://github.com/math-proof/lemma}}
\date{}

\begin{document}
\maketitle

\begin{abstract}
Ferrari's method reduces a general quartic to a depressed quartic and a cubic resolvent; Cardano then solves that cubic by cube roots.
Over an arbitrary field the identities are classical, and Mathlib's archived Ferrari theorem leaves the resolvent root and the square roots as extra data.
What is not classical is a \emph{total} map \(\C^5\to\C^4\) in \emph{principal} radicals: Mathlib's \(\sqrt{\,\cdot\,}\) and \(\sqrt[3]{\,\cdot\,}\) (Lean \texttt{Root.cubic}, notation \(\sqrt[3]{\,}\)) fail the product law, so Cardano's pairing \(3AB=-p\) does not hold for two independent principal cube roots.

We state and machine-check such a map in the \texttt{lemma} library.
Every auxiliary quantity is a principal square or cube root; the missing cube root of unity is the radical \(-\tfrac12+i\sqrt{3}/2\) (the exponential \(e^{2\pi i/3}\) is an abbreviation only); the branch is the single power \(\omega^{k}\) with
\[
k=\Bigl\lceil\frac{3\Arg(-p/3)}{2\pi}-\tfrac12\Bigr\rceil
-\Bigl\lceil\frac{3\Arg(AB)}{2\pi}-\tfrac12\Bigr\rceil,
\]
not a case split on \(k\bmod 3\).
The proof is checked in Lean~4 (updated 2026-08-30) and was first stated interactively in SymPy (2018-11-29).
Interactive visualizations and source files are public on \url{http://www.lemma.cn/} and GitHub~\cite{lemma_repo}.
\end{abstract}

\section{Introduction}
\label{sec:intro}

Let \(a,b,c,d,e,x\in\C\) with \(a\neq 0\) and
\[
a x^4+b x^3+c x^2+d x+e=0.
\]
Ferrari (16th century) translates to a depressed quartic \(z^4+\alpha z^2+\beta z+\gamma=0\), produces a cubic resolvent, and factors the quartic into two quadratics once a resolvent root is known.
Cardano solves that cubic by cube roots.
Over \(\C\), Mathlib's \(\sqrt{z}\) and \(\sqrt[3]{z}\) are defined by the principal logarithm (argument in \((-\pi,\pi]\)).
They satisfy
\[
(\sqrt{z})^2=z,\qquad \bigl(\sqrt[3]{z}\bigr)^3=z
\]
for every \(z\), but they do \emph{not} satisfy
\[
\sqrt{z}\sqrt{w}=\sqrt{zw},\qquad \sqrt[3]{z}\,\sqrt[3]{w}=\sqrt[3]{zw}
\]
in general.
Any fully radical formula over \(\C\) must therefore name the missing cube root of unity explicitly and select a consistent Cardano pair.

This paper records the public lemma
\texttt{Complex.Imp\_OrOrSEqS\_Sub.Imp\_OrOrSEqS\_SubDiv\_2.of.Eq0Add\_Mul\_Pow\_4.Ne\_0}~\cite{lemma_quartic_lean,lemma_quartic_sympy}:
a conjunction of two implications (\(\beta=0\) and \(\beta\neq 0\)), each a four-fold disjunction of nested principal radicals.
The cube-root branch is \(y=A\omega^{k}+B\), not three residue classes of an integer modulo~3.
Cube roots in the statement are Lean \(\sqrt[3]{\,}\) (\texttt{Root.cubic}), not a bare \texttt{cpow} written as \(z^{(3:\C)^{-1}}\).

The same module is visualized on lemma.cn in SymPy~\cite{lemma_quartic_sympy} and Lean~4~\cite{lemma_quartic_lean}.
Source files are \texttt{Ne\_0.lean} (branch \texttt{main}) and \texttt{Ne\_0.py} (branch \texttt{master})~\cite{lemma_repo}.

\section{Related work}
\label{sec:related}

\textbf{Mathlib~4.}
Zhu's archive formalization~\cite{zhu2024mathlibquartic} (PR~\#18290; Wiedijk items 37 and 46~\cite{wiedijk100}) proves Ferrari over a field \(K\) of characteristic not two.
Square roots appear as hypotheses \(t^2=\cdots\), \(v^2=\cdots\), and the resolvent root \(u\) is an extra datum.
The comment on \texttt{quartic\_eq\_zero\_iff} states that an explicit Cardano expression for \(u\) ``would be too long.''
The theorem is not a map \(\C^5\to\C^4\) in principal radicals, and it is not specific to \(\C\)~\cite{mathlib}.

\textbf{Dyson--Ahrens--Emmenegger}~\cite{dyson2022polynomials}
work in Lean~3 in a field equipped with \emph{some} square-root and cube-root operations.
The functions are axioms, not Mathlib's principal branches; the product identity for cube roots is deliberately not assumed.
Instantiating the typeclass with \(\C\) does not prove the principal-branch formula.

\textbf{Coq}
(the development compared in that paper, \S7.1) is specific to \(\C\) and uses De~Moivre \(n\)th roots.
That is a complex-domain radical solution, but it is not a Lean~4 lemma and does not use Mathlib's principal \(\sqrt{\,\cdot\,}\) / \texttt{Root.cubic} together with an \(\Arg\)/\(\lceil\cdot\rceil\) selector.

\textbf{Isabelle/HOL} \texttt{Cubic\_Quartic} and the HOL Light entry on Wiedijk's list are algebraic in the same sense as Mathlib: roots are assumed, not constructed by principal functions.

This library already has a companion cubic lemma~\cite{lemma_cubic_lean}
that solves \(ax^3+bx^2+cx+d=0\) over \(\C\) by principal radicals and the same integer \(k\), writing the three Cardano roots as \(A\omega^{k}+B\), \(A\omega^{k-1}+B\omega\), and \(A\omega^{k+1}+B\,\overline{\omega}\).
The quartic contribution is to finish Ferrari on that calculus: the resolvent is expanded by Cardano, one legal root \(y=A\omega^{k}+B\) is selected, and the four quartic roots are nested principal square roots of that \(y\).

\section{The formula}
\label{sec:formula}

Write \(a'=b/a\), \(b'=c/a\), \(c'=d/a\), \(d'=e/a\), and
\begin{align*}
\alpha&=b'-3a'^2/8,\\
\beta&=a'^3/8+c'-a'b'/2,\\
\gamma&=a'^2 b'/16+d'-3a'^4/256-a'c'/4.
\end{align*}
Then \(z=x+a'/4\) satisfies \(z^4+\alpha z^2+\beta z+\gamma=0\).

\paragraph{Biquadratic regime \(\beta=0\).}
With \(\Delta=\alpha^2-4\gamma\),
\[
x=\pm\sqrt{\bigl(\pm\sqrt{\Delta}-\alpha\bigr)/2}-a'/4
\]
(four independent choices of the two signs).

\paragraph{Ferrari regime \(\beta\neq 0\).}
Form Cardano quantities of the resolvent,
\begin{align*}
p&=-4\gamma-\alpha^2/3,\\
q&=-2\alpha^3/27+8\alpha\gamma/3-\beta^2,\\
\delta&=4p^3/27+q^2,\\
A&=\sqrt[3]{(-q+\sqrt{\delta})/2},\qquad
B=\sqrt[3]{(-q-\sqrt{\delta})/2},
\end{align*}
and
\[
k=\Bigl\lceil\frac{3\Arg(-p/3)}{2\pi}-\tfrac12\Bigr\rceil
-\Bigl\lceil\frac{3\Arg(AB)}{2\pi}-\tfrac12\Bigr\rceil.
\]
The abbreviation \(\omega=e^{2\pi i/3}\) (Lean \texttt{(I * (2 * $\pi$ / 3)).exp}) is for simplicity only: a complex exponential is \emph{not} a radical.
The same primitive cube root of unity is the principal-radical rectangular form
\[
\omega=-\tfrac12+i\,\tfrac{\sqrt{3}}{2}
\]
(Lean \texttt{(-(1 / 2 : $\mathbb{R}$)) + I * ($\sqrt{3}$ / 2 : $\mathbb{R}$)}), and \(\overline{\omega}=-\tfrac12-i\,\tfrac{\sqrt{3}}{2}\).
The identities \(e^{2\pi i/3}=-\tfrac12+i\sqrt{3}/2\) and its conjugate are machine-checked~\cite{lemma_omega_rect,lemma_omega_conj}.
Set
\[
y=A\,\omega^{k}+B,\qquad
y_0=-2\alpha/3+y,\qquad
y_1=4\alpha/3+y.
\]
Then
\begin{align*}
x&=\frac{\sqrt{2\beta/\sqrt{y_0}-y_1}-\sqrt{y_0}}{2}-\frac{a'}{4}
\quad\text{or}\quad
\frac{-\sqrt{2\beta/\sqrt{y_0}-y_1}-\sqrt{y_0}}{2}-\frac{a'}{4}\\
&\quad\text{or}\quad
\frac{\sqrt{-2\beta/\sqrt{y_0}-y_1}+\sqrt{y_0}}{2}-\frac{a'}{4}
\quad\text{or}\quad
\frac{-\sqrt{-2\beta/\sqrt{y_0}-y_1}+\sqrt{y_0}}{2}-\frac{a'}{4}.
\end{align*}

\begin{theorem}[Principal-radical quartic~\cite{lemma_quartic_lean}]
\label{thm:quartic}
Assume \(a\neq 0\) and \(ax^4+bx^3+cx^2+dx+e=0\).
With the quantities above:
\begin{enumerate}[noitemsep]
  \item if \(\beta=0\), then \(x\) equals one of the four biquadratic expressions;
  \item if \(\beta\neq 0\), then \(x\) equals one of the four Ferrari expressions with \(y=A\omega^{k}+B\).
\end{enumerate}
The cube-root branch is the single factor \(\omega^{k}\), not a case split on \(k\bmod 3\).
\end{theorem}

The statement does not claim that every combination of nested radicals is a root for a wrong branch; it claims that a genuine root \(x\) equals one of the four expressions once that power of \(\omega\) is selected.

\section{Figures}
\label{sec:figures}

\paragraph{Figure~0 --- public visualizations.}
The SymPy page~\cite{lemma_quartic_sympy} displays two implications (\(\beta=0\), \(\beta\neq 0\)) as an interactive theorem document (created 2018-11-29).
The Lean~4 page~\cite{lemma_quartic_lean} displays \texttt{@[main] private lemma main} with \(\{x,a,b,c,d,e:\C\}\), the \texttt{let} chain \(a',b',\alpha,\beta,\gamma,p,q,\delta,A,B,k,\omega,y,y_0,y_1\) (\(A,B\) via \(\sqrt[3]{\,}\)), and the conjunction of two radical disjunctions.

\paragraph{Figure~1 --- reduction.}
Figure~\ref{fig:reduction} is the reduction used in the main theorem: divide by the leading coefficient, depress by \(z=x+a'/4\), then split on \(\beta=0\) versus \(\beta\neq 0\).

\begin{figure}[ht]
\centering
\footnotesize
\[
\begin{array}{c}
a x^{4}+b x^{3}+c x^{2}+d x+e=0,\quad a\neq 0 \\[0.6em]
\Big\downarrow\ \scriptstyle a'=\frac{b}{a},\; b'=\frac{c}{a},\; c'=\frac{d}{a},\; d'=\frac{e}{a} \\[0.6em]
x^{4}+a' x^{3}+b' x^{2}+c' x+d'=0 \\[0.6em]
\Big\downarrow\ \scriptstyle z=x+\frac{a'}{4} \\[0.6em]
z^{4}+\alpha z^{2}+\beta z+\gamma=0 \\[0.25em]
\scriptstyle
\alpha=b'-\frac{3a'^{2}}{8},\quad
\beta=\frac{a'^{3}}{8}+c'-\frac{a'b'}{2},\quad
\gamma=\frac{a'^{2}b'}{16}+d'-\frac{3a'^{4}}{256}-\frac{a'c'}{4}
\\[1.0em]
\begin{array}{c@{\qquad}c}
\beta=0 & \beta\neq 0 \\[0.4em]
\Big\downarrow & \Big\downarrow \\[0.4em]
\Delta=\alpha^{2}-4\gamma &
y=A\,\omega^{k}+B
\\[1.6em]
x=\pm\sqrt{\dfrac{\pm\sqrt{\Delta}-\alpha}{2}}-\dfrac{a'}{4}
&
\begin{array}{c}
y_{0}=-\dfrac{2\alpha}{3}+y,\quad y_{1}=\dfrac{4\alpha}{3}+y \\[0.45em]
x=\dfrac{\pm\sqrt{\pm\dfrac{2\beta}{\sqrt{y_{0}}}-y_{1}}\pm\sqrt{y_{0}}}{2}-\dfrac{a'}{4}
\end{array}
\end{array}
\\[1.0em]
x=z-\dfrac{a'}{4}
\end{array}
\]
\caption{Reduction of the general quartic and the two radical regimes.}
\label{fig:reduction}
\end{figure}

\paragraph{Figure~2 --- lemma dependencies.}
Figure~\ref{fig:deps} lists major imported lemmas of Theorem~\ref{thm:quartic}.
Private helpers in the same file are drawn in parentheses.
Links point to the Lean~4 pages; the SymPy counterpart of each module is the same path under \url{http://www.lemma.cn/py/?module=}.

\begin{figure}[ht]
\small
\begin{itemize}[leftmargin=1.2em,itemsep=0.15em]
  \item \href{http://www.lemma.cn/lean/?module=Complex.Imp_OrOrSEqS_Sub.Imp_OrOrSEqS_SubDiv_2.of.Eq0Add_Mul_Pow_4.Ne_0}{\texttt{Imp\_OrOrSEqS\_Sub.\ldots Ne\_0}} --- main theorem
  \begin{itemize}[leftmargin=1.1em,itemsep=0.1em]
    \item \href{http://www.lemma.cn/lean/?module=Complex.OrOrSEqS.of.Eq0Add_Pow_4}{\texttt{OrOrSEqS.of.Eq0Add\_Pow\_4}} --- case \(\beta=0\)
    \begin{itemize}[leftmargin=1.1em,itemsep=0.05em]
      \item \href{http://www.lemma.cn/lean/?module=Complex.OrEqS_Div.of.Eq0Add_Mul_Square.Ne_0}{\texttt{OrEqS\_Div.of.Eq0Add\_Mul\_Square.Ne\_0}} --- quadratic in \(x^2\)
      \item \href{http://www.lemma.cn/lean/?module=Complex.Or_Eq_NegSqrt.of.EqSquare}{\texttt{Or\_Eq\_NegSqrt.of.EqSquare}} --- \(x^2=c\Rightarrow x=\pm\sqrt{c}\)
      \begin{itemize}[leftmargin=1.1em,itemsep=0.05em]
        \item \href{http://www.lemma.cn/lean/?module=Complex.EqSquareSqrt}{\texttt{EqSquareSqrt}} --- \((\sqrt{z})^2=z\)
        \item \href{http://www.lemma.cn/lean/?module=Real.OrEqS.of.Square}{\texttt{Real.OrEqS.of.Square}}
      \end{itemize}
    \end{itemize}
    \item \href{http://www.lemma.cn/lean/?module=Complex.OrOrSEqS_Div_2.of.Eq0Add_Pow_4.Ne_0}{\texttt{OrOrSEqS\_Div\_2.of.Eq0Add\_Pow\_4.Ne\_0}} --- case \(\beta\neq 0\)
    \begin{itemize}[leftmargin=1.1em,itemsep=0.05em]
      \item (\texttt{hbranch}) --- given a resolvent root \(y_0\), factor into two quadratics
      \item \href{http://www.lemma.cn/lean/?module=Complex.Eq_Mul_Pow_SubCeilS.of.Pow_3}{\texttt{Eq\_Mul\_Pow\_SubCeilS.of.Pow\_3}} --- \(A^3=B^3\Rightarrow A=B\,\omega^{k}\)
      \begin{itemize}[leftmargin=1.1em,itemsep=0.05em]
        \item \href{http://www.lemma.cn/lean/?module=Complex.ArgPow.eq.SubMul_Arg}{\texttt{ArgPow.eq.SubMul\_Arg}}
        \item \href{http://www.lemma.cn/lean/?module=Complex.Eq_MulNorm_ExpMulIArg}{\texttt{Eq\_MulNorm\_ExpMulIArg}}
        \item \href{http://www.lemma.cn/lean/?module=Complex.ArgExpMulI.eq.Sub_Mul_Ceil}{\texttt{ArgExpMulI.eq.Sub\_Mul\_Ceil}}
        \item \href{http://www.lemma.cn/lean/?module=Int.Floor.eq.NegCeilNeg}{\texttt{Int.Floor.eq.NegCeilNeg}}
      \end{itemize}
    \end{itemize}
  \end{itemize}
\end{itemize}
\caption{Import graph for the Lean proof of Theorem~\ref{thm:quartic}.}
\label{fig:deps}
\end{figure}

Related lemmas in the same family, \emph{not} imported by the main theorem: the depressed statement
\href{http://www.lemma.cn/lean/?module=Complex.Imp_OrOrSEqS.Imp_OrOrSEqS_Div_2.of.Eq0Add_Pow_4}{\texttt{Imp\_OrOrSEqS.Imp\_OrOrSEqS\_Div\_2.of.Eq0Add\_Pow\_4}}
and the monic offset
\href{http://www.lemma.cn/lean/?module=Complex.Imp_OrOrSEqS_Sub.Imp_OrOrSEqS_SubDiv_2.of.Eq0Add_Pow_4}{\texttt{Imp\_OrOrSEqS\_Sub.Imp\_OrOrSEqS\_SubDiv\_2.of.Eq0Add\_Pow\_4}}.
The cubic formula~\cite{lemma_cubic_lean} uses the same \(k\) selector; the quartic proof calls \texttt{Eq\_Mul\_Pow\_SubCeilS.of.Pow\_3} directly, because Ferrari needs one resolvent root \(A\omega^{k}+B\) rather than three Cardano combinations.

\paragraph{Figure~3 --- branch integer \(k\).}
From \(A^3=B^3\) one has \(A=B\,\omega^{k}\) with
\[
k=\Bigl\lceil\frac{3\Arg A}{2\pi}-\tfrac12\Bigr\rceil
-\Bigl\lceil\frac{3\Arg B}{2\pi}-\tfrac12\Bigr\rceil.
\]
Applied to \(AB\) and \(-p/3\), this is the same \(k\) as in Theorem~\ref{thm:quartic}, and \(\omega^{-k}\) restores \(AB=(-p/3)\,\omega^{-k}\).
Multiplying only \(A\) by \(\omega^{k}\) produces a legal Ferrari parameter \(y=A\omega^{k}+B\).
Because \(\omega^{3}=1\), the three residue classes of \(k\) modulo~3 are already encoded in that single power: \(\omega^{k}\in\{1,\omega,\overline{\omega}\}\).

\section{Proof outline}
\label{sec:proof}

\subsection{Principal radicals}
\label{sec:radicals}

Throughout, \(\sqrt{\,\cdot\,}\) is Mathlib \texttt{Complex.sqrt} and \(\sqrt[3]{\,\cdot\,}\) is Lean \texttt{Root.cubic}.
Both are total functions \(\C\to\C\).
The identity \((\sqrt{z})^2=z\) is \texttt{Complex.EqSquareSqrt}.
If \(x^2=c\), then \(x=\sqrt{c}\) or \(x=-\sqrt{c}\) by \texttt{Complex.Or\_Eq\_NegSqrt.of.EqSquare}.
Cube roots of unity are the radicals
\[
\omega=-\tfrac12+i\,\tfrac{\sqrt{3}}{2},\qquad
\overline{\omega}=-\tfrac12-i\,\tfrac{\sqrt{3}}{2},\qquad
\omega^{2}=\overline{\omega},\qquad
\omega^{3}=1.
\]
The polar form \(z=\|z\|e^{i\Arg z}\) is \texttt{Complex.Eq\_MulNorm\_ExpMulIArg}.
Reduction of \(\Arg(e^{ix})\) into \((-\pi,\pi]\) by a ceiling is \texttt{Complex.ArgExpMulI.eq.Sub\_Mul\_Ceil}, which yields \texttt{Complex.ArgPow.eq.SubMul\_Arg} and then
\texttt{Complex.Eq\_Mul\_Pow\_SubCeilS.of.Pow\_3}:
if \(A^3=B^3\), then \(A=B\,\omega^{k}\) for the integer \(k\) of Figure~3.

\subsection{Depression}
\label{sec:depress}

Division by \(a\neq 0\) produces the monic equation \(x^4+a'x^3+b'x^2+c'x+d'=0\).
The translation \(z=x+a'/4\) cancels the cubic term and yields \(z^4+\alpha z^2+\beta z+\gamma=0\) with \(\alpha,\beta,\gamma\) as in \S\ref{sec:formula}.
The main file performs this expansion by \texttt{ring}; the four nested-radical expressions for \(z\) are then shifted back by \(-a'/4\).

\subsection{Biquadratic case \(\beta=0\)}
\label{sec:biquad}

The depressed equation is \(z^4+\alpha z^2+\gamma=0\), a quadratic in \(z^2\).
Lemma \texttt{Complex.OrOrSEqS.of.Eq0Add\_Pow\_4} assembles the four principal-radical roots.
It first solves
\[
z^2=\frac{-\alpha\pm\sqrt{\alpha^2-4\gamma}}{2}=\frac{\pm\sqrt{\Delta}-\alpha}{2},\qquad \Delta=\alpha^2-4\gamma,
\]
then splits each \(z^2=c\) by \texttt{Or\_Eq\_NegSqrt}.
Translating \(x=z-a'/4\) is the first conjunct of Theorem~\ref{thm:quartic}.

\subsection{Ferrari factorization}
\label{sec:ferrari}

Let \(y_0\in\C\) be any root of the cubic resolvent
\[
y_0^3+2\alpha y_0^2+(\alpha^2-4\gamma)y_0-\beta^2=0,
\]
and set \(y_1=y_0+2\alpha\).
A private helper proves \(y_0\neq 0\) (else \(\beta=0\)), that \((\sqrt{y_0})^2=y_0\), and that with \(Y=(y_0+\alpha)/2\) one has
\[
(z^2+Y)^2-\bigl(\sqrt{y_0}\,z-\beta/(2\sqrt{y_0})\bigr)^2
=z^4+\alpha z^2+\beta z+\gamma.
\]
The left-hand side is a product of two quadratics.
Each is solved by completing the square and \texttt{Or\_Eq\_NegSqrt}, producing
\begin{align*}
z&=\frac{\sqrt{2\beta/\sqrt{y_0}-y_1}-\sqrt{y_0}}{2},&
z&=\frac{-\sqrt{2\beta/\sqrt{y_0}-y_1}-\sqrt{y_0}}{2},\\
z&=\frac{\sqrt{-2\beta/\sqrt{y_0}-y_1}+\sqrt{y_0}}{2},&
z&=\frac{-\sqrt{-2\beta/\sqrt{y_0}-y_1}+\sqrt{y_0}}{2}.
\end{align*}
This step still treats \(y_0\) as given.
The rest of the proof constructs \(y_0\) from principal radicals.

\subsection{Cardano data of the resolvent}
\label{sec:cardano}

The theorem statement writes the scaled Cardano pair directly:
\begin{align*}
p&=-4\gamma-\alpha^2/3,\\
q&=-2\alpha^3/27+8\alpha\gamma/3-\beta^2,\\
\delta&=4p^3/27+q^2,\\
A&=\sqrt[3]{(-q+\sqrt{\delta})/2},\qquad
B=\sqrt[3]{(-q-\sqrt{\delta})/2}.
\end{align*}
The proof also depresses the resolvent by
\[
a_r=-\alpha/2,\qquad b_r=-\gamma,\qquad c_r=-\beta^2/8+\alpha\gamma/2,
\]
and the unscaled Cardano parameters
\begin{align*}
p_c&=b_r-a_r^2/3,\qquad
q_c=2a_r^3/27-a_r b_r/3+c_r,\qquad
\delta_c=4p_c^3/27+q_c^2,\\
A_c&=\sqrt[3]{(-q_c+\sqrt{\delta_c})/2},\qquad
B_c=\sqrt[3]{(-q_c-\sqrt{\delta_c})/2}.
\end{align*}
A direct calculation gives \(q=8q_c\), \(\delta=64\delta_c\), and \(\sqrt{\delta}=8\sqrt{\delta_c}\) (the last because \(64>0\) on the principal branch).
Positive real scaling of cube roots yields \(A=2A_c\) and \(B=2B_c\), and \(\Arg(AB)=\Arg(A_c B_c)\).

\subsection{Restoring \(3(A\omega^{k})B=-p\)}
\label{sec:product}

Because \((AB)^3=(-p/3)^3\), lemma \texttt{Eq\_Mul\_Pow\_SubCeilS.of.Pow\_3} supplies
\[
AB=(-p/3)\,\omega^{-k}.
\]
Multiplying by \(\omega^{k}\) restores the product in one step: \(AB\,\omega^{k}=-p/3\).
Since \((\omega^{k})^3=1\),
\[
(A\omega^{k})^3+B^3=A^3+B^3=-q,\qquad 3(A\omega^{k})B=-p,
\]
so \(y=A\omega^{k}+B\) satisfies \(y^3+py+q=0\).
There is no case split on \(k\bmod 3\): the three possible rotations \(1\), \(\omega\), \(\overline{\omega}\) are the values of the single power \(\omega^{k}\).
Shifting \(y_0=-2\alpha/3+y\) and \(y_1=4\alpha/3+y\) recovers a resolvent root, and the four nested square roots apply.

Unlike the cubic formula in this library, only \(A\) is rotated.
A quartic needs one legal resolvent parameter, not three Cardano roots of the original unknown.

\section{Discussion and limitations}
\label{sec:discussion}

The statement is a conjunction of two implications, not an if-and-only-if that would list roots for a mismatched branch, and not a four-fold split on \(k\bmod 3\).
Each inner disjunction is exhaustive for that case.
Principal square roots in the Ferrari step may vanish or repeat when the quartic has multiple roots; the disjunction still holds.

The lemma does not construct roots in a general field of characteristic zero, does not avoid \(\Arg\) and \(\lceil\cdot\rceil\), and does not prove that a randomly chosen pairing of principal cube roots is a resolvent root.
It does not replace Mathlib's abstract Ferrari theorem~\cite{zhu2024mathlibquartic}: that theorem remains the right statement when one only knows that some square roots and some resolvent root exist in \(K\).
The present lemma is the complementary statement on \(\C\): every coefficient tuple with \(a\neq 0\) is sent to four explicit principal-radical expressions, with the cube-root branch named by the power \(\omega^{k}\).

\section{Conclusion}
\label{sec:conclusion}

We have presented a machine-checked, fully expanded Ferrari--Cardano formula for the general quartic over \(\C\), using principal square roots, principal cube roots \(\sqrt[3]{\,}\), the radical cube root of unity \(-\tfrac12+i\sqrt{3}/2\), and a single power \(\omega^{k}\) assembled from \(\Arg\) and \(\lceil\cdot\rceil\).
The result is public on lemma.cn and GitHub~\cite{lemma_repo}, with proofs checked in Lean~4 and an earlier SymPy formulation.

\bibliographystyle{plain}
\bibliography{refs}

\end{document}
